%% ─────────────────────────────────────────────────────────────────────────────
%% CV — Talvin Ackbaraly (Version française)
%% Serif (XCharter), one restrained accent colour, thin rules, bullets, no icons
%% Compile: latexmk -pdf cv-talvin-ackbaraly-fr.tex
%% ─────────────────────────────────────────────────────────────────────────────
\documentclass[10.5pt, a4paper]{extarticle}

\usepackage[T1]{fontenc}
\usepackage[utf8]{inputenc}
\usepackage[french]{babel}
\usepackage[top=1cm, bottom=1cm, left=1.6cm, right=1.6cm]{geometry}
\usepackage{XCharter}
\usepackage[letterspace=80]{microtype}
\usepackage{titlesec}
\usepackage{enumitem}
\usepackage{tabularx}
\usepackage{ragged2e}
\usepackage{xcolor}
\usepackage[scaled=0.94,noupquote]{inconsolata}
\usepackage{tikz}
\usepackage{varwidth}
\usepackage{graphicx}
\usepackage[hidelinks]{hyperref}

\hypersetup{pdftitle={CV — Talvin Ackbaraly}, pdfauthor={Talvin Ackbaraly}}

%% ── Colour system: one accent, everything else neutral ───────────────────────
\definecolor{accent}{HTML}{2563A0}  % THE accent: name, sections, links, bullets
\definecolor{ink}{HTML}{1A1A1A}     % body text (warm near-black, not pure 000)
\definecolor{soft}{HTML}{6E6E6E}    % secondary information (context, tech)
\definecolor{hair}{HTML}{C9C9C9}    % section rules
\definecolor{chipbg}{HTML}{EAF0F7}  % tag pill background
\definecolor{chipfg}{HTML}{2F5B87}  % tag pill text

\hypersetup{colorlinks=true, urlcolor=accent, linkcolor=accent}
\AtBeginDocument{\color{ink}}

\pagestyle{empty}
\setlength{\parindent}{0pt}
\setlength{\parskip}{0pt}
\linespread{1.02}
\RaggedRight
\hyphenpenalty=10000
\exhyphenpenalty=10000
\emergencystretch=1em
\microtypesetup{expansion=false}

%% Section: letter-spaced small caps over a hairline
\titleformat{\section}
  {\large\scshape\lsstyle\color{accent}}{}{0pt}{}
  [\vspace{-5pt}{\color{hair}\rule{\textwidth}{0.4pt}}]
\titlespacing*{\section}{0pt}{13pt}{5pt}

\frenchsetup{StandardItemLabels=true}

%% Tag pill: \chip{CUDA}
\newcommand{\chip}[1]{%
  \tikz[baseline=(c.base)]{\node[inner xsep=4.5pt, inner ysep=2pt, rounded corners=2.5pt,
    fill=chipbg, text=chipfg, font=\footnotesize] (c) {#1};}\hspace{1pt}%
}
%% \entry{Title}{Date}{Subtitle}{Tags}
\newcommand{\entry}[4]{%
  \textbf{#1} \hfill #2\\
  \textit{#3} \hfill {#4}\par
}
%% \project{Name}{Context}{Tags}
\newcommand{\project}[3]{%
  \textbf{#1}{\color{soft}\enspace—\enspace\textit{#2}} \hfill {#3}\par
}
\newcommand{\gap}{\par\vspace{6pt}}


%% NB : en tête de cellule « p », {\color{...}} décale d'une ligne — toujours \textcolor.

%% ── 1. Formation : entrées empilées, sans tags ───────────────────────────────
\newcommand{\edu}[4]{%
  \textbf{#1}\hfill #2\par\vspace{1.5pt}
  \textcolor{accent}{\textit{#3}}\par\vspace{1.5pt}
  {\small\color{soft}#4}\par
}

%% ── 2. Projets : prose (titre en tête) / tags ────────────────────────────────
%% Chaque ligne calcule sa propre largeur de colonne de tags : on mesure la
%% largeur naturelle des pastilles, on la plafonne, et la prose prend le reste.
\newlength{\projgapb}\setlength{\projgapb}{0.4cm}
\newlength{\projmaxtag}\setlength{\projmaxtag}{3.35cm}
\newlength{\projtagw}\newlength{\projprosew}
\newsavebox{\projtagbox}
%% #1 titre (en tête de prose, sans point final), #2 prose, #3 pastilles.
%% varwidth compose les pastilles jusqu'à \projmaxtag mais rétrécit à la largeur
%% réellement occupée : chaque ligne garde donc sa propre gouttière.
\newcommand{\proj}[3]{%
  \sbox\projtagbox{\begin{varwidth}[t]{\projmaxtag}\raggedleft #3\end{varwidth}}%
  \setlength{\projtagw}{\wd\projtagbox}%
  \setlength{\projprosew}{\dimexpr\textwidth-\projgapb-\projtagw\relax}%
  \noindent
  \parbox[t]{\projprosew}{\RaggedRight \textbf{#1.} #2}%
  \hspace{\projgapb}%
  \usebox\projtagbox
  \par
}
\newcommand{\projgap}{\vspace{10pt}}

%% ── 3. Expérience : rôle d'abord, prose pleine largeur ───────────────────────
%% #1 rôle, #2 date, #3 employeur, #4 pastilles, #5 prose.
%% La date occupe le coin haut droit — c'est le premier point de fixation du
%% regard ; les pastilles descendent sur la ligne de l'employeur.
\newcommand{\job}[5]{%
  \textbf{#1}\hfill #2\par\vspace{1.5pt}
  \textcolor{accent}{#3}\hfill {#4}\par\vspace{4pt}
  #5\par
}

%% ── 4. Compétences : libellé + pastilles ─────────────────────────────────────
\newlength{\skilllab}\setlength{\skilllab}{4.95cm}
\newcommand{\skillrow}[2]{%
  \noindent\parbox[t]{\skilllab}{\textbf{#1}}%
  \parbox[t]{\dimexpr\textwidth-\skilllab\relax}{\RaggedRight #2}\par\vspace{3.5pt}
}


%% ── En-tête : photo ronde en haut à droite ───────────────────────────────────
%% Photo : assets/photo.jpg ; sinon un emplacement vide s'affiche.
\newlength{\photosize}\setlength{\photosize}{3.4cm}
\IfFileExists{assets/photo.jpg}{\def\photofile{assets/photo.jpg}}{%
  \IfFileExists{../../assets/photo.jpg}{\def\photofile{../../assets/photo.jpg}}{}}
\newcommand{\headerphoto}{%
  \ifdefined\photofile
    \tikz{\clip (0,0) circle (0.5\photosize);
      \node at (0,0) {\includegraphics[width=\photosize]{\photofile}};
      \draw[hair, line width=0.6pt] (0,0) circle (0.5\photosize);}%
  \else
    \tikz{\draw[hair, fill=chipbg, line width=0.6pt] (0,0) circle (0.5\photosize);
      \node[text=chipfg, font=\footnotesize] at (0,0) {Photo};}%
  \fi
}

\begin{document}

\noindent
\begin{minipage}[c]{\dimexpr\textwidth-\photosize-0.3cm\relax}
  \RaggedRight
  {\color{accent}\fontsize{25}{30}\selectfont\scshape\textls[40]{Talvin Ackbaraly}}\\[6pt]
  \textbf{Stage de fin d'études — Vision par ordinateur \& IA · 6 mois dès février 2027}\\[5pt]
  {\small
  \href{mailto:talvin.ackbaraly@icloud.com}{talvin.ackbaraly@icloud.com} \enspace·\enspace
  \href{tel:+33769389008}{07 69 38 90 08} \enspace·\enspace
  \href{https://talvin-ackbaraly.com}{talvin-ackbaraly.com}\\[2pt]
  \href{https://github.com/talvinckb}{github.com/talvinckb} \enspace·\enspace
  \href{https://www.linkedin.com/in/talvin-ackbaraly}{linkedin.com/in/talvin-ackbaraly}}\\[7pt]
  Étudiant en dernière année à EPITA, majeure Image. Après un parcours en développement
  full-stack et ingénierie logicielle, je souhaite mettre en pratique ma spécialisation en vision
  par ordinateur et IA en entreprise. Basé en Île-de-France.
\end{minipage}\hfill
\begin{minipage}[c]{\photosize}
  \headerphoto
\end{minipage}

\section{Expérience professionnelle}

\job{Projet de fin d'études — Vision par ordinateur \textnormal{\color{soft}(équipe de 4)}}{mai 2026 -- en cours}
    {EPITA $\times$ Bibliothèque nationale de France (Gallica)}{\chip{PyTorch}\chip{YOLO26}\chip{ConvNeXt}\chip{Label Studio}}
    {Détection (YOLO26) et classification hiérarchique (ConvNeXt) d'illustrations patrimoniales
     dans les documents numérisés de Gallica, via l'API IIIF. En charge de la détection : entraînement sur
     un dataset de 10\,000 images annotées par la BnF sur Label Studio, benchmark de
     modèles, réglage des hyperparamètres, pour \textbf{94,7\,\%~de~mAP50}
     (85,2\,\% en mAP50-95) sur le jeu de validation.}
\vspace{11pt}
\job{Développeur Full-Stack \textnormal{\color{soft}(freelance)}}{sept. 2026 -- en cours}
    {Auto-entrepreneur}{\chip{Next.js}\chip{TypeScript}\chip{NestJS}\chip{PostgreSQL}}
    {Lancement d'une activité de développement web en freelance. Mission en cours : application de
     gestion de stock pour un client professionnel, prise en charge de bout en bout, du cahier des charges
     au déploiement.}
\vspace{11pt}
\job{Développeur Full-Stack \& DevOps \textnormal{\color{soft}(stage)}}{sept. 2025 -- janv. 2026}
    {Studiocanal}{\chip{Java}\chip{Spring Boot}\chip{Angular}\chip{AWS}}
    {Développement de fonctionnalités full-stack (back-end Java/Spring Boot, front-end Angular) pour une
     plateforme interne de gestion des droits, contrats et catalogue, utilisée par 100+ collaborateurs. Mise en place du pipeline CI/CD
     GitLab d'un nouveau module et intervention sur l'infrastructure AWS (EC2, RDS, IAM), en équipe Agile.}

\section{Projets}

\proj{Segmentation d'images avec U-Net}{U-Net (encodeur-décodeur, skip connections) implémenté from scratch en
  PyTorch, sans modèle pré-entraîné : segmentation premier plan / arrière-plan sur COCO, avec boucle d'entraînement
  et évaluation sur données non vues.}
  {\chip{PyTorch}\chip{U-Net}\chip{COCO}}
\projgap
\proj{Prédiction de fonction pulmonaire}{Estimation de la capacité vitale forcée à partir de scanners CT :
  segmentation pulmonaire et XGBoost, livré en application web React conteneurisée.}
  {\chip{Python}\chip{OpenCV}\chip{XGBoost}\chip{Docker}}
\projgap
\proj{Lecture de plaques d'immatriculation}{Détection classique (filtres, morphologie, HOG) et Random Forest :
  \textbf{86\,\% de précision}. Benchmark Python vs C++, workflow GitLab complet.}
  {\chip{OpenCV}\chip{C++}\chip{Random Forest}}
\projgap
\proj{Suivi de tumeurs cérébrales}{Recalage 3D d'IRM (rigide, affine, B-spline), segmentation de gliome
  et quantification de l'évolution tumorale (+61,8\,\% de volume), avec visualisation 3D.}
  {\chip{ITK}\chip{VTK}\chip{PyQt6}}


\section{Formation}

\edu{EPITA — École d'ingénieurs en informatique}{sept. 2022 -- juil. 2027}
    {Diplôme d'ingénieur, majeure Image — Paris}
    {Traitement \& synthèse d'images, vision par ordinateur, deep learning.}
\vspace{7pt}
\edu{Université du Québec à Chicoutimi (UQAC)}{janv. -- mai 2024}
    {Semestre d'échange académique — Québec, Canada}
    {Informatique générale : UML, programmation orientée objet, cybersécurité, réalité virtuelle.}


\section{Compétences}

\skillrow{Vision par ordinateur \& IA}{\chip{Traitement d'images}\chip{Deep learning}\chip{PyTorch}\chip{OpenCV}\chip{YOLO}\chip{U-Net}\chip{ConvNeXt}\chip{XGBoost}}
\skillrow{Programmation}{\chip{C}\chip{C++}\chip{Python}\chip{Java}\chip{TypeScript}\chip{SQL}}
\skillrow{Développement \& DevOps}{\chip{Spring Boot}\chip{Angular}\chip{React/Next.js}\chip{NestJS}\chip{Docker}\chip{GitLab CI}\chip{AWS}\chip{PostgreSQL}}
\skillrow{Langues}{Français (langue maternelle), anglais (professionnel, TOEIC 845)}

\end{document}
